% SageFs Formal Verification — Conference Paper
% Lean Squad automated formal verification for WillEhrendreich/SageFs
%
% ⚠️ LaTeX not available in CI sandbox — PDF not compiled.
%    Compile with: cd formal-verification/paper && latexmk -pdf paper.tex

\documentclass[sigconf,11pt]{acmart}

\usepackage{listings}
\usepackage{booktabs}
\usepackage{amsmath,amssymb}
\usepackage{xcolor}
\usepackage{graphicx}
\usepackage{microtype}

%% Lean 4 listing style
\lstdefinelanguage{Lean4}{
  keywords={theorem, def, lemma, example, structure, inductive, where,
            import, open, namespace, end, sorry, by, have, let, in,
            match, with, if, then, else, do, return, partial, Type,
            Prop, fun, forall, exists, And, Or, Not, True, False},
  sensitive=true,
  morecomment=[l]{--},
  morecomment=[n]{/-}{-/},
  morestring=[b]",
}
\lstset{
  language=Lean4,
  basicstyle=\ttfamily\small,
  keywordstyle=\bfseries,
  commentstyle=\itshape\color{gray!70},
  stringstyle=\color{red!70!black},
  breaklines=true,
  frame=single,
  numbers=left,
  numberstyle=\tiny\color{gray},
  xleftmargin=2em,
  framexleftmargin=1.5em,
}

\title{Formal Verification of an F\# Live Development Environment with Lean~4}
\subtitle{177 Theorems, 11 Files, and a Compositional Evaluation-Gate Proof}

\author{Lean Squad}
\affiliation{%
  \institution{Automated Formal Verification Agent}
  \country{GitHub Actions}}
\email{lean-squad@github-actions.bot}

\begin{abstract}
We report on a formal verification effort for SageFs, a live F\# development
environment consisting of a REPL daemon, editor integrations, and a state-driven
evaluation pipeline. Using Lean~4 with the standard library (no Mathlib, due to CI
network constraints), we produced 11 specification files containing 177 theorems
proved with zero \texttt{sorry}. The verification covers nine core modules
individually---including a circular ring buffer, railway-oriented error combinators,
retry/restart policy models, a tool-gating affordance layer, a computation-expression
evaluation trace, a file-watching state machine, a session phase projector, and a
theme configuration overlay---and two cross-cutting files that reason about the system
compositionally. The central system-level result is the evaluation-gate theorem
\texttt{send\_fsharp\_code\_iff\_ready\_phase}, which proves that code submission is
available \emph{if and only if} the session phase is \texttt{Active·Idle}, unifying
two independent modules into a single formal statement. A \texttt{validTransition}
inductive relation models all eight session state-machine transitions and establishes
safety invariants including ``evaluation failure never leads directly to the Faulted
state.'' No implementation bugs were found; rather, the proofs confirmed key design
invariants and revealed structural properties such as the unreachability of
\texttt{SessionState.Uninitialized} via the normal lifecycle. Runnable correspondence
tests (50 cases) validate the ring-buffer model against the F\# implementation.
\end{abstract}

%%
%% Keywords
\keywords{formal verification, Lean 4, F\#, functional programming, state machines,
  proof assistant, developer tools, session lifecycle}

\maketitle

%% ============================================================
\section{Introduction}
%% ============================================================

Developer tooling is increasingly complex: a modern live development environment
orchestrates a long-running language interpreter, file-system watchers, editor
integrations, an evaluation pipeline, and a state machine governing session lifecycle.
Correctness properties that seem obvious by inspection---``code evaluation is only
available when the session is ready''---span multiple subsystems and are easy to
violate during refactoring.

SageFs~\cite{sagefs} is an F\# live development environment providing a persistent
F\# Interactive (FSI) session, Hot Reload, a terminal UI, Raylib GUI client, and
editor integrations for VS~Code, Visual Studio, and Neovim. Its core logic in
\texttt{SageFs.Core} includes:
\begin{itemize}
  \item A \texttt{RingBuffer<'T>} circular buffer for storing evaluation history.
  \item Railway-oriented combinators (\texttt{ResultEx}) for error threading.
  \item Retry and restart policy models with exponential backoff.
  \item An \texttt{Affordances} module gating tool availability by session state.
  \item An \texttt{EvalPipeline} computation-expression for traced execution.
  \item \texttt{HotReloadState} file-watching state machine.
  \item A session lifecycle model (\texttt{AppState}) with five phases.
  \item A \texttt{Theme} configuration overlay combinator.
\end{itemize}

We apply Lean~4 formal verification to this codebase with three goals: (1) confirm
correctness of key data-structure invariants and policy decisions; (2) establish
system-level composition theorems connecting independent modules; and (3) detect any
implementation defects through proof failure or counterexample.

\paragraph{Contributions.}
\begin{enumerate}
  \item 177 Lean~4 theorems covering nine SageFs core modules, proved with zero
    \texttt{sorry}.
  \item The first cross-file composition theorem in the project:
    \texttt{send\_fsharp\_code\_iff\_ready\_phase} unifies the session lifecycle and
    tool-affordance subsystems.
  \item A \texttt{validTransition} inductive relation formalising all eight
    session state-machine transitions with safety invariants.
  \item 50 runnable correspondence tests validating the ring-buffer Lean model against
    the F\# implementation.
  \item A methodology for applying Lean~4 stdlib-only verification to an effectful
    functional codebase, including explicit documentation of all modelling
    approximations.
\end{enumerate}

%% ============================================================
\section{Background}
%% ============================================================

\subsection{The SageFs Codebase}

SageFs is implemented in F\# targeting \texttt{net10.0}. The project uses a daemon
architecture: a long-running FSI session managed by a mailbox processor
(\texttt{AppState.fs}), which transitions through five phases: \texttt{Uninitialized},
\texttt{Initializing}, \texttt{Active(Idle)}, \texttt{Active(Evaluating)}, and
\texttt{Faulted}. Editor communication uses an MCP server; the terminal UI uses
SageTUI's Elm Architecture. The core library (\texttt{SageFs.Core}) contains the
modules targeted for formal verification.

\subsection{Lean 4 and the Standard Library}

Lean~4~\cite{lean4} is a dependently typed theorem prover and functional programming
language. It supports both tactic-based and term-mode proofs, with a rich standard
library including decidable propositions, natural number arithmetic (\texttt{Nat}),
lists, arrays, and \texttt{Except} for error handling.

We use Lean~4 stdlib only (version 4.30.0-rc2, toolchain
\texttt{leanprover/lean4:v4.30.0-rc2}). Mathlib~\cite{mathlib4} is unavailable due to
CI network constraints (the Azure blob cache endpoint is blocked). This constrains our
automation to \texttt{decide}, \texttt{omega}, \texttt{simp}, and \texttt{cases}/\texttt{induction},
with no access to Mathlib's extended tactic library.

\subsection{Related Work}

Formal verification of functional programming languages and their runtimes has a rich
history. The CompCert~\cite{compcert} project verified a C compiler; several efforts
have targeted Haskell~\cite{liquidhaskell} and OCaml~\cite{ocaml-fv} core libraries.
Lean~4 itself has been used to verify mathematical results in Mathlib and to target
systems software~\cite{lean4-systems}. To our knowledge, no prior work has targeted an
F\# live development environment, nor the specific class of session state-machine and
tool-affordance properties we verify here.

%% ============================================================
\section{Methodology}
%% ============================================================

\subsection{Target Selection}

We surveyed \texttt{SageFs.Core} for FV-amenable targets: pure or nearly-pure
functions, data structure invariants, and policy-enforcement logic. We applied
the following criteria:
\begin{itemize}
  \item \textbf{Pure functional core}: the target should have a clear input/output
    mapping separable from I/O and mutation.
  \item \textbf{Testable invariants}: existing tests, documentation, or naming
    conventions signal which properties must hold.
  \item \textbf{Security or safety relevance}: access-control policies and session
    state machines have high bug-catching potential.
  \item \textbf{Non-trivial structure}: targets that require reasoning about
    modular arithmetic, list membership, or state transitions are more valuable
    than trivially decidable ones.
\end{itemize}

Nine primary targets were identified (Table~\ref{tab:targets}), plus two
cross-cutting files added after the individual modules matured.

\begin{table}[h]
\centering
\small
\begin{tabular}{llp{4.5cm}}
\toprule
Module & Theorems & Primary property \\
\midrule
\texttt{RingBuffer} & 20 & \texttt{WellFormed} invariant, push/get correctness \\
\texttt{ResultEx} & 17 & Monad/functor laws, sequence/partition \\
\texttt{RetryPolicy} & 13 & Decision correctness, delay monotonicity \\
\texttt{RestartPolicy} & 9 & Exponential backoff correctness \\
\texttt{Affordances} & 19 & Tool-gating access-control policy \\
\texttt{EvalPipeline} & 17 & Computation-expression trace structure \\
\texttt{HotReloadState} & 23 & Watch/unwatch/toggle invariants \\
\texttt{SessionLifecycle} & 16 & Phase projection, state reachability \\
\texttt{Theme} & 20 & Override identity and idempotency \\
\texttt{Composition} & 12 & Cross-file evaluation gate \\
\texttt{PhaseTransition} & 11 & Transition relation + safety invariants \\
\midrule
\textbf{Total} & \textbf{177} & — \\
\bottomrule
\end{tabular}
\caption{Verified modules and theorem counts.}
\label{tab:targets}
\end{table}

\subsection{Specification Strategy}

For each target we followed a two-phase approach:
\begin{enumerate}
  \item \textbf{Informal spec extraction}: we wrote a Markdown document
    (\texttt{formal-verification/specs/<name>\_informal.md}) describing preconditions,
    postconditions, invariants, edge cases, and open questions inferred from the code.
  \item \textbf{Lean formalisation}: we translated the informal spec into Lean~4 type
    definitions, function stubs, and \texttt{theorem} declarations, initially guarded
    by \texttt{sorry}, then progressively proved.
\end{enumerate}

\subsection{Modelling Choices}

Each F\# module is modelled as a pure functional Lean definition. The key
approximations, documented in \texttt{CORRESPONDENCE.md}, are:

\begin{itemize}
  \item \textbf{Immutability}: F\# mutable records and arrays become immutable Lean
    structures and \texttt{Array} values. The observable input/output semantics are
    preserved.
  \item \textbf{Integer types}: F\# \texttt{int} (signed 64-bit) becomes Lean \texttt{Nat}
    (non-negative). The \texttt{WellFormed} predicate ensures non-negativity; overflow
    at $2^{63}$ is not modelled.
  \item \textbf{I/O and effects}: file system access, FSI session communication, and
    timing are omitted entirely. Proofs apply to the pure functional core.
  \item \textbf{Concurrency}: the mailbox processor and async transitions in
    \texttt{AppState.fs} are modelled as sequential transitions. Concurrent
    \texttt{cancelEval} racing with evaluation is not captured.
  \item \textbf{Parametric abstraction}: the session phase type \texttt{Phase α} is
    kept polymorphic; the type parameter is not constrained.
\end{itemize}

\subsection{Proof Approach}

We applied Lean~4 tactics in order of increasing complexity:
\begin{enumerate}
  \item \textbf{Decidable propositions}: \texttt{decide} closes goals over finite,
    concrete values (string lengths, specific constants).
  \item \textbf{Arithmetic}: \texttt{omega} for natural-number arithmetic and
    \texttt{simp\_arith} for combined simplification.
  \item \textbf{Structural}: \texttt{simp only [...]} with explicit lemma lists to
    unfold definitions; \texttt{cases} and \texttt{induction} for sum types and
    inductive propositions.
  \item \textbf{Case splits}: \texttt{by\_cases (key' == key)} for Boolean predicates;
    \texttt{constructor}/\texttt{intro}/\texttt{apply}/\texttt{exact} for logical
    connectives.
  \item \textbf{Modular arithmetic}: manual lemmas (\texttt{index\_after\_push},
    \texttt{newHead\_ne\_index}) for the ring-buffer head-position proofs.
\end{enumerate}

\subsection{Model Validation}

To validate correspondence between the Lean ring-buffer model and the F\# source,
we wrote a runnable test harness in F\# script (\texttt{run.fsx}) under
\texttt{formal-verification/tests/ringbuffer/}. The harness exercises 50 test cases
covering normal behaviour, edge cases (eviction, clear, single-element buffer),
and property-based tests (20 sequential pushes). All 50 tests pass.

Continuous integration is provided by \texttt{.github/workflows/lean-ci.yml}, which
runs \texttt{lake build} on every PR touching \texttt{formal-verification/lean/**},
and \texttt{fv-correspondence-tests.yml}, which runs the F\# correspondence tests on
changes to source or test files.

%% ============================================================
\section{Results}
%% ============================================================

\subsection{Proof Inventory}

Table~\ref{tab:theorems} lists the key proved theorems by module and property.
All 177 theorems are proved with zero \texttt{sorry}; \texttt{lake build} passes
with Lean 4.30.0-rc2.

\begin{table}[h]
\centering
\small
\begin{tabular}{p{4.2cm}lp{3.8cm}}
\toprule
Theorem & File & Property \\
\midrule
\texttt{push\_wellFormed} & RingBuffer & Invariant preserved by push \\
\texttt{push\_aging} & RingBuffer & Modular-arithmetic aging \\
\texttt{toList\_length} & RingBuffer & List length = count \\
\texttt{resBind\_assoc} & ResultEx & Bind associativity (monad law) \\
\texttt{resSequence\_length} & ResultEx & Sequence preserves length \\
\texttt{baseDelay\_mono} & RetryPolicy & Delay grows monotonically \\
\texttt{rtDecide\_giveup\_exhausted} & RetryPolicy & Give up when retries gone \\
\texttt{send\_fsharp\_code\_iff\_ready} & Affordances & Eval gate ↔ Ready state \\
\texttt{hard\_reset\_denied\_when\_eval} & Affordances & Reset denied during eval \\
\texttt{two\_step\_first\_fails\_one\_stage} & EvalPipeline & Trace on early failure \\
\texttt{toggle\_involution} & HotReloadState & Double toggle = identity \\
\texttt{watch\_idempotent} & HotReloadState & Watch is idempotent \\
\texttt{toState\_never\_uninitialized} & SessionLifecycle & Unreachability finding \\
\texttt{ready\_iff\_active\_idle} & SessionLifecycle & State classification \\
\texttt{withOverrides\_empty\_id} & Theme & Empty override = identity \\
\texttt{withOverrides\_idempotent\_single} & Theme & Single override idempotency \\
\texttt{stateToSessionState\_bijective} & Composition & Bridge isomorphism \\
\texttt{send\_fsharp\_code\_iff\_ready\_phase} & Composition & End-to-end eval gate \\
\texttt{evaluating\_cannot\_send\_code} & Composition & Reentrancy guard \\
\texttt{tools\_always\_available} & Composition & Tool always accessible \\
\texttt{faulted\_only\_recovers\_to\_init} & PhaseTransition & Recovery path constraint \\
\texttt{eval\_cannot\_fault\_directly} & PhaseTransition & Eval failure → Ready \\
\texttt{uninitialized\_unreachable\_as\_target} & PhaseTransition & Transition safety \\
\bottomrule
\end{tabular}
\caption{Selected key theorems from all 177 proved.}
\label{tab:theorems}
\end{table}

\subsection{Composition Theorems}

The two cross-cutting files represent a qualitative step beyond individual-module
verification.

\paragraph{Composition.lean.}
\texttt{Composition.lean} imports both \texttt{FVSquad.SessionLifecycle} and
\texttt{FVSquad.Affordances}. It proves that the two \texttt{State}/\texttt{SessionState}
inductives (independently defined in each module) are isomorphic via
\texttt{stateToSessionState}, and then leverages this bridge to state system-level
properties:

\begin{lstlisting}[language=Lean4,caption={Central evaluation gate theorem.}]
theorem send_fsharp_code_iff_ready_phase
    {α : Type} (p : Phase α) :
    checkToolAvailability
      (stateToSessionState (toState p))
      "send_fsharp_code" ↔
    ∃ s, p = Phase.Active s Activity.Idle := by
  cases p <;> simp [toState, stateToSessionState,
                    checkToolAvailability,
                    availableTools]
\end{lstlisting}

This theorem is the first formal statement in the project connecting the lifecycle
and affordance subsystems. It asserts that code submission is available \emph{exactly
when} the session is in the \texttt{Active·Idle} state.

\paragraph{PhaseTransition.lean.}
\texttt{PhaseTransition.lean} defines a \texttt{validTransition} inductive
relation with eight constructors derived from the eval actor's message handlers in
\texttt{AppState.fs} (Table~\ref{tab:transitions}).

\begin{table}[h]
\centering
\small
\begin{tabular}{lll}
\toprule
Constructor & From & To \\
\midrule
\texttt{initToInit} & Initializing & Initializing \\
\texttt{initToReady} & Initializing & Active(Idle) \\
\texttt{initToFaulted} & Initializing & Faulted \\
\texttt{readyToEval} & Active(Idle) & Active(Evaluating) \\
\texttt{evalToReady} & Active(Evaluating) & Active(Idle) \\
\texttt{readyToInit} & Active(Idle) & Initializing \\
\texttt{evalToInit} & Active(Evaluating) & Initializing \\
\texttt{faultedToInit} & Faulted & Initializing \\
\bottomrule
\end{tabular}
\caption{\texttt{validTransition} cases and their F\# counterparts.}
\label{tab:transitions}
\end{table}

The key safety invariant \texttt{eval\_cannot\_fault\_directly} proves:

\begin{lstlisting}[language=Lean4,caption={Evaluation never leads directly to Faulted.}]
theorem eval_cannot_fault_directly
    {α : Type} (s : Session α)
    (t : Phase α) :
    validTransition
      (Phase.Active s Activity.Evaluating) t →
    t ≠ Phase.Faulted := by
  intro h; cases h <;> simp
\end{lstlisting}

This confirms the F\# design: the \texttt{EvalFinished(Error ex)} handler sets
\texttt{sessionState'' = SessionState.Ready}, not Faulted. Faulted is only
reachable via explicit reset operations.

\subsection{Structural Finding: Uninitialized is Unreachable}

A notable structural finding from \texttt{SessionLifecycle.lean} is that
\texttt{SessionState.Uninitialized} (a 5-case state type used by the Affordances
layer) is unreachable via the \texttt{toState} projection from the lifecycle's
\texttt{Phase} type:

\begin{lstlisting}[language=Lean4,caption={Unreachability of Uninitialized.}]
theorem toState_never_uninitialized
    {α : Type} (p : Phase α) :
    toState p ≠ State.Uninitialized := by
  cases p <;> simp [toState]
\end{lstlisting}

This confirms that the external API cannot observe \texttt{Uninitialized} from a
well-typed session phase. The \texttt{validTransition} relation corroborates this
at the transition level: \texttt{uninitialized\_unreachable\_as\_target} proves
no valid transition leads to \texttt{Uninitialized}.

\subsection{Bugs and Findings}

No implementation bugs were found in any of the nine primary targets. All stated
properties hold for the Lean functional models.

One \textbf{model inaccuracy} was discovered and corrected during development (D3
in CORRESPONDENCE.md): an early version of \texttt{rbClear} reset \texttt{total}
to 0, diverging from the F\# \texttt{clear} which preserves \texttt{Total}. This
was caught during correspondence review and corrected before any theorem relied on
it.

\subsection{Coverage Assessment}

Table~\ref{tab:coverage} summarises which aspects of each module are covered.

\begin{table}[h]
\centering
\small
\begin{tabular}{lp{3.5cm}p{3cm}}
\toprule
Module & Covered & Omitted \\
\midrule
RingBuffer & Invariants, push/get/clear semantics & Mutable arrays, overflow \\
ResultEx & Monad/functor laws, sequence/partition & Exception types \\
RetryPolicy & Decision logic, delay formula & Wall-clock time \\
RestartPolicy & Backoff correctness & Window expiry \\
Affordances & Tool-gating policy & Role-based auth \\
EvalPipeline & Trace structure & FSI I/O, timing \\
HotReloadState & Watch/unwatch/toggle & File system events \\
SessionLifecycle & Phase projection, reachability & Async transitions \\
Theme & Override semantics & File-based loading \\
Composition & Eval gate, bridge isomorphism & Side effects \\
PhaseTransition & All 8 transitions, safety invariants & Concurrency \\
\bottomrule
\end{tabular}
\caption{Coverage summary per module.}
\label{tab:coverage}
\end{table}

The primary gap is concurrency: the mailbox processor and async cancel/reset
racing are not modelled. The proofs establish correctness of the sequential
state-machine abstraction; additional work would be needed to handle interleaving.

%% ============================================================
\section{Discussion}
%% ============================================================

\subsection{Proof Utility}

The highest-value theorems are those connecting independently specified subsystems:
\texttt{send\_fsharp\_code\_iff\_ready\_phase} and the \texttt{validTransition}
safety invariants. These would catch regressions introduced by refactoring that
changes the session phase logic or the tool-gating policy in isolation. The
lower-value theorems---monad laws, string-length checks---are correct but add
little beyond what unit tests already provide.

A concern noted in CRITIQUE.md is that the 18 \texttt{withOverrides\_*\_preserves\_*}
theorems in \texttt{Theme.lean} each verify a pair of fields independently rather
than the whole record. A stronger field-completeness theorem would subsume all 18
but requires decidable equality on field names.

\subsection{Automation and Effort}

The project was produced entirely by the Lean Squad automated agent across
approximately 11 runs. Human effort was limited to reviewing and merging PRs.
The most time-consuming proofs were the ring-buffer modular arithmetic lemmas
(requiring manual helper lemmas) and the Theme field-isolation theorems (requiring
18 separate case-split instances). The composition theorems, by contrast, were
surprisingly simple once the bridge function was defined: most closed with
\texttt{cases p <\&> simp [...]}.

\subsection{Limitations}

\begin{itemize}
  \item \textbf{No Mathlib}: the absence of Mathlib limits automation. Several proofs
    that would be one-line \texttt{simp} calls with Mathlib lemmas required manual
    helper lemmas.
  \item \textbf{Concurrency not modelled}: the async session lifecycle is abstracted
    to a sequential transition relation.
  \item \textbf{No extraction}: we do not extract Lean code to F\# or generate
    certified implementations. The proofs are post-hoc verification of a
    hand-written model.
  \item \textbf{Correspondence by code review}: for most modules, the correspondence
    between the Lean model and the F\# source is validated by manual inspection and
    the ring-buffer correspondence tests. Automated extraction (e.g., via Aeneas for
    a Rust codebase) is not applicable here.
\end{itemize}

\subsection{Lessons Learned}

\begin{enumerate}
  \item \textbf{Start with invariants, not algorithms}: the most tractable early
    targets were data-structure invariants (\texttt{WellFormed}, \texttt{toList\_length})
    rather than full functional correctness.
  \item \textbf{Compose late}: individual-module proofs are easier to write in
    isolation; composition theorems should be added once the modules are stable.
  \item \textbf{Document approximations immediately}: the D3 divergence (rbClear total
    reset) was caught early because the correspondence review happened alongside proof
    writing. Deferring this review would have allowed incorrect theorems to accumulate.
  \item \textbf{Stdlib-only is viable}: without Mathlib, proofs require more manual
    lemmas, but the approach scales to 177 theorems with manageable effort per theorem.
\end{enumerate}

%% ============================================================
\section{Conclusion}
%% ============================================================

We have presented a formal verification effort for SageFs, an F\# live development
environment, using Lean~4 with the standard library. The project produced 177 proved
theorems across 11 files, with zero \texttt{sorry} and CI-checked builds. The central
contribution is the system-level evaluation-gate theorem
\texttt{send\_fsharp\_code\_iff\_ready\_phase}, which unifies the session lifecycle
and tool-affordance subsystems in a single formal statement. The
\texttt{validTransition} relation establishes safety invariants for the session
state machine, and 50 runnable correspondence tests validate the ring-buffer model
against the F\# implementation.

No implementation bugs were found, but the process revealed a model inaccuracy
(D3, ring-buffer clear semantics) and the structural unreachability of
\texttt{SessionState.Uninitialized} via the normal lifecycle. Future work includes
modelling the async session transitions, adding correspondence tests for
\texttt{HotReloadState} and \texttt{Affordances}, and strengthening the Theme
field-isolation coverage.

The project demonstrates that Lean~4 formal verification is practical for a
real-world F\# developer tool, even without Mathlib and with a fully automated
verification agent.

%% ============================================================
\bibliographystyle{ACM-Reference-Format}
\bibliography{paper}

\end{document}
